\section{0x31\ 数学\ 多项式}
\label{sec:0x31}

\subsection{拉格朗日插值}

\Highlight{多项式的插值}：给出 $n + 1$ 个点 $(x_1, y_1), (x_2, y_2), \dots, (x_{n + 1}, y_{n + 1})$，\textbf{其中横坐标（模意义下）两两不同}。求一个\textbf{至多 $n$ 次}的多项式，使得这 $n + 1$ 个点都在 $f(x)$ 上。

\Highlight{拉格朗日插值}：

\[
f(x) = \sum\limits_{i = 1}^{n + 1} y_i \prod\limits_{j \neq i} \frac{x - x_j}{x_i - x_j}
\]

可以发现 $f(x)$ 是一个\textbf{至多 $n$ 次}的多项式，并且可以证明\textbf{拉格朗日插值是唯一的插值多项式}。

朴素实现的时间复杂度为 $\mathcal{O}(n^2)$，可以使用多项式优化到 $\mathcal{O}(n \log^2 n)$（详见多项式快速插值）。

\begin{quoteblock}
构造方法：考虑构造 $n + 1$ 个函数，\textbf{其中第 $i$ 个函数 $f_i(x)$ 的图像经过 $(x_i, y_i)$，对于其他的 $j \neq i$ 经过 $(x_j, 0)$}。则我们要求的 $f(x)$ 即为这 $n + 1$ 个函数的和 $\sum_{i = 1}^{n + 1} f_i(x)$。

尝试构造 $f_i(x) = a_i\prod_{j \neq i}(x - x_j)$，其中 $a_i$ 是待定的系数。此时对于 $j \neq i$ 均有 $f_i(x_j) = 0$。代入 $f_i(x_i) = y_i$ 即可解得 $a_i = \frac{y_i}{\prod_{j \neq i}(x_i - x_j)}$。

于是原式得证。
\end{quoteblock}

\InputCode{cpp}{contents/0x30 数学/code/0x31 拉格朗日插值.cpp}

\Highlight{拉格朗日插值求系数}：相当于是直接将整个多项式求出来。运用了\textbf{暴力除以单项式}的技巧。

\begin{itemize}
\item 记系数部分 $a_i = \frac{y_i}{\prod_{j \neq i} (x_i - x_j)}$。
\item 记多项式 $g(x) = \prod_{j} (x - x_j)$，记多项式 $h_i(x) = \prod_{j \neq i} (x - x_j)$。
\item 先 $\mathcal{O}(n^2)$ 暴力多项式乘法求出 $g(x)$。\textbf{考虑由 $g(x)$ 得到 $h_i(x)$，相当于是多项式 $g(x)$ 除以一个单项式 $(x - x_i)$}，仍然可以 $\mathcal{O}(n)$ 暴力除以单项式。具体地，由于 $(x - x_i)h_i(x) = g(x)$，则对于 $g(x)$ 的 $j$ ($j \geq 1$) 次项系数，有 $[x^{j - 1}]h_i - x_i[x^j]h_i = [x^j]g$。于是

\[
\begin{gathered}
[x^n]h_i = [x^{n + 1}]g \\
[x^j]h_i = [x^{j + 1}]g + x_i[x^{j + 1}]h_i \\
\end{gathered}
\]

\item 求出了 $a_i$ 与 $h_i(x)$ 之后，计算 $f(x) = \sum_{i = 1}^{n + 1}a_ih_i(x)$。时间复杂度 $\mathcal{O}(n^2)$。
\end{itemize}

\InputCode{cpp}{contents/0x30 数学/code/0x31 拉格朗日插值求系数.cpp}

\Highlight{\textbf{连续点值}的拉格朗日插值}：给出 $n + 1$ 个连续点值 $f(1), f(2), \dots, f(n + 1)$，代入拉格朗日插值公式，得

\[
f(x) = \sum\limits_{i = 1}^{n + 1}y_i \prod\limits_{j \neq i}\frac{x - j}{i - j}
\]

现在考虑给定一个自变量 $x$，求出因变量 $f(x)$。考虑右边的乘积项怎么求

\begin{itemize}
\item 分子：记 $\mathrm{pre}_i = \prod_{1 \leq j \leq i} (x - j)$，$\mathrm{suf}_i = \prod_{i \leq j \leq n + 1} (x - j)$，则分子为 $\mathrm{pre}_{i - 1} \times \mathrm{suf}_{i + 1}$。
\item 分母：$j < i$ 的部分\textbf{是 $i - 1$ 到 $1$ 乘起来}。$j > i$ 的部分\textbf{是 $-1$ 到 $-(n + 1 - i)$ 乘起来}。故分母为 $(i - 1)! (n + 1 - i)! (-1)^{n + 1 - i}$。
\end{itemize}

故最后的式子为

\[
f(x) = \sum\limits_{i = 1}^{n + 1}y_i \frac{\mathrm{pre}_{i - 1} \times \mathrm{suf}_{i + 1}}{(i - 1)!(n + 1 - i)!(-1)^{n + 1 - i}}
\]

求出单个点值的时间复杂度为 $\mathcal{O}(n)$。

\begin{quoteblock}
若给出 $n + 1$ 个连续点值 $f(h), f(h + 1), \dots, f(h + n)$，分母的求法是完全不变的，只需将分子改成 $x - x_j$ 相乘即可。

或者可以考虑将坐标系向左平移 $h - 1$。
\end{quoteblock}

\InputCode{cpp}{contents/0x30 数学/code/0x31 连续点值的拉格朗日插值.cpp}

\subsection{FFT}

\Highlight{单位根}：$n$ 次单位根，即为 $x^n = 1$ 的所有复数解。单位圆上辐角为 $0, \frac{2\pi}{n}, \frac{4\pi}{n}, \dots, \frac{2(n - 1)\pi}{n}$ 的复数都是单位根，共 $n$ 个。这些单位根依次表示为

\[
\omega_n^0, \omega_n^1, \omega_n^2, \dots, \omega_n^{n - 1}
\]
\Highlight{单位根的简单性质}：欧拉公式：$e^{ix} = \cos(x) + i \sin(x)$。这说明 $x \in \mathbb{R}$ 时，$e^{ix}$ 可以描述一个单位圆上辐角为 $x$ 的复数。

\begin{itemize}
\item $\omega_n^1 = e^{\frac{2\pi i}{n}} = \cos(\frac{2\pi}{n}) + i \sin(\frac{2\pi}{n})$。
\item $\omega_n^i \times \omega_n^j = \omega_n^{i + j}$。
\item $\left(\omega_n^i\right)^j = \omega_n^{ij}$。
\item $\omega_{2n}^{2k} = \omega_n^k$。
\item $\omega_{2n}^k = -\omega_{2n}^{k + n}$。
\end{itemize}

\Highlight{补充界}：在进行 FFT 之前，需要对其最高位补 $0$，将其补成一个最高次数为 $2^P - 1$ 的多项式。方便后续的奇偶性分组。

\subsubsection{DFT}

DFT 的思想：“系数表示法” 转 “点值表示法”。

已知系数 $a_0, \dots, a_{n - 1}$，取自变量 $\omega_n^0, \omega_n^1, \dots, \omega_n^{n - 1}$，计算点值 $F(\omega_n^0), F(\omega_n^1), \dots, F(\omega_n^{n - 1})$。

记 $n = 2^P$，考虑一个 $2^P - 1$ 次的多项式 $F(x)$

\[
F(x) = a_0 + a_1x + \dots + a_{n - 1}x^{n - 1}
\]
将每一项按照次数的奇偶性，划分成两部分（由于 $n$ 是 $2$ 的若干次幂，每次分出来的奇偶部分的次数一定相等）。

\[
F(x) = \left( a_0 + a_2x^2 + \dots + a_{n - 2}x^{n - 2} \right) + \left( a_1x + a_3x^3 + \dots + a_{n - 1}x^{n - 1} \right)
\]
记

\[
\begin{gathered}
F_L(x) = a_0 + a_2x + \dots + a_{n - 2}x^{n / 2 - 1} \\
F_R(x) = a_1 + a_3x + \dots + a_{n - 1}x^{n / 2 - 1}
\end{gathered}
\]
则

\[
F(x) = F_L(x^2)+ xF_R(x^2)
\]
考虑计算因变量 $F(\omega_n^0), F(\omega_n^1), \dots, F(\omega_n^{n - 1})$。

\begin{itemize}
\item 对于 $k < n / 2$，代入 $\omega_n^k$ 得

\[
\begin{aligned}
F(\omega_n^k) & = F_L\left((\omega_n^k)^2\right) + \omega_n^k F_R\left((\omega_n^k)^2\right) \\
& = F_L(\omega_{n / 2}^k) + \omega_n^k F_R(\omega_{n / 2}^k)
\end{aligned}
\]

\item 对于 $k < n / 2$，代入 $\omega_n^{k + n/2}$ 得

\[
\begin{aligned}
F(\omega_n^{k + n / 2}) & = F_L\left((\omega_n^{k + n / 2})^2\right) + \omega_n^{k + n / 2}F_R\left((\omega_n^{k + n / 2})^2\right) \\
& = F_L(\omega_n^{2k + n}) + \omega_n^{k + n / 2} F_R(\omega_n^{2k + n}) \\
& = F_L(\omega_{n / 2}^k) - \omega_{n}^{k} F_R(\omega_{n / 2}^k)
\end{aligned}
\]
\end{itemize}

注意到 $F(\omega_n^k), F(\omega_n^{k + n / 2})$ 等于 $F_L(\omega_{n / 2}^k) \pm \omega_n^kF_R(\omega_{n / 2}^k)$。

如果我们知道 $F_L(x), F_R(x)$ 分别在 $\omega_{n / 2}^0, \omega_{n / 2}^1, \dots, \omega_{n / 2}^{n / 2 - 1}$ 处的点值表示，就可以 $\mathcal{O}(n)$ 求出 $F(x)$ 在 $\omega_n^0, \omega_n^1, \dots, \omega_n^{n - 1}$ 处的点值表示。

这是一个分治的过程，不断分治直到仅剩一个项即可。

时间复杂度 $\mathcal{O}(n \log n)$。

\subsubsection{IDFT}

\textbf{结论：将 DFT 中的 $\omega_n^1$ 换成 $\omega_n^{-1}$，做完以后除以 $n$ 即可}。

IDFT 的思想：“点值表示法” 转 “系数表示法”。

已知点值 $F(\omega_n^0), F(\omega_n^1), \dots, F(\omega_n^{n - 1})$，求出系数 $a_0, \dots, a_{n - 1}$。

记 $n = 2^P$，对于一个 $2^{P} - 1$ 的多项式 $F(x)$，设其经过 DFT 之后，得到的点值序列 $g_0, g_1, \dots, g_{n - 1}$，可以证明：

\[
g_k = \sum\limits_{i = 0}^{n - 1}a_i \left( \omega_n^k \right)^i \iff a_k = \frac{1}{n} \sum\limits_{i = 0}^{n - 1}g_i \left( \omega_n^{-k} \right)^i
\]
证明：将 $g_i$ 代入该式子，得

\[
\begin{aligned}
  & \frac{1}{n} \sum\limits_{i = 0}^{n - 1}\sum\limits_{j = 0}^{n - 1}(\omega_n^{-k})^i(\omega_n^i)^j a_j \\
= & \sum\limits_{j = 0}^{n - 1} a_j \left(\frac{1}{n}\sum\limits_{i = 0}^{n - 1}\omega_{n}^{i(j - k)}\right)
\end{aligned}
\]
后面括号里的和式，即为单位根反演的经典结论 $[n \mid k] = \frac{1}{n}\sum_{i = 0}^{n - 1}\omega_{n}^{ik}$。具体地，分类讨论：

\begin{itemize}
\item 当 $j = k$ 时，和式为 $\frac{1}{n}\sum_{i = 0}^{n - 1}\omega_n^0 = 1$。
\item 当 $j \neq k$ 时，和式为

\[
\frac{1}{n} \sum_{i = 0}^{n - 1} \omega_{n}^{i(j - k)} = \frac{1}{n}\sum_{i = 0}^{n - 1} \left(\omega_n^{j - k} \right)^i = \frac{1}{n} \frac{1 - (\omega_n^{j - k})^n}{1 - \omega_n^{j - k}} = 0
\]
\end{itemize}

故 $a_k = \frac{1}{n} \sum_{i = 0}^{n - 1}g_i \left( \omega_n^{-k} \right)^i$ 成立。

\Separator{}

进一步，观察 IDFT 的式子，相当于将 DFT 中的 $\omega_n^1$ 换成 $\omega_n^{-1}$，做完以后再除以 $n$。

再进一步，发现 $F(\omega_n^{-i}) = F(\omega_n^{n - i})$，于是 IDFT 可以看作将 DFT 后得到的点值进行翻转（注意 $F(\omega_n^0)$ 这一项不参与翻转），然后再除以 $n$。

套用 DFT 的板子即可。

\subsubsection{*DFT 与 IDFT 的矩阵角度}

对于 DFT 前的系数表示法 $a_0, a_1, \dots, a_{n - 1}$ 以及 DFT 后的点值表示法 $g_0, g_1, \dots, g_{n - 1}$。可以写成一个\textbf{系数向量}左乘一个转移矩阵得到\textbf{点值向量}的式子：

\[
\begin{bmatrix}
g_0 \\
g_1 \\
\vdots \\
g_{n - 1}
\end{bmatrix}
=
\begin{bmatrix}
1 & 1 & 1 & 1 & \dots & 1 \\
1 & \omega_n^1 & \omega_n^2 & \omega_n^3 & \dots & \omega_n^{n - 1} \\
1 & \omega_n^2 & \omega_n^4 & \omega_n^6 & \dots & \omega_n^{2(n - 1)} \\
1 & \omega_n^3 & \omega_n^6 & \omega_n^9 & \dots & \omega_n^{3(n - 1)} \\
\vdots & \vdots & \vdots & \vdots & \ddots & \vdots \\
1 & \omega_n^{n - 1} & \omega_n^{2(n - 1)} & \omega_n^{3(n - 1)} & \dots & \omega_n^{(n - 1)(n - 1)}
\end{bmatrix}
\begin{bmatrix}
a_0 \\
a_1 \\
\vdots \\
a_{n - 1}
\end{bmatrix}
\]
其中转移矩阵（下标从 $0$ 开始）第 $i$ 行第 $j$ 列的系数即为自变量 $\omega_n^i$ 的 $j$ 次方 $\left(\omega_{n}^i\right)^j$。

如何从点值向量得到系数向量？只需让上式等号两边都左乘一个转移矩阵的逆即可。\textbf{结论：DFT 矩阵的逆，为各个系数取倒数，再除以矩阵大小 $n$}。故有

\[
\begin{bmatrix}
a_0 \\
a_1 \\
\vdots \\
a_{n - 1}
\end{bmatrix}
=
\frac{1}{n}
\begin{bmatrix}
1 & 1 & 1 & 1 & \dots & 1 \\
1 & \omega_n^{-1} & \omega_n^{-2} & \omega_n^{-3} & \dots & \omega_n^{-(n - 1)} \\
1 & \omega_n^{-2} & \omega_n^{-4} & \omega_n^{-6} & \dots & \omega_n^{-2(n - 1)} \\
1 & \omega_n^{-3} & \omega_n^{-6} & \omega_n^{-9} & \dots & \omega_n^{-3(n - 1)} \\
\vdots & \vdots & \vdots & \vdots & \ddots & \vdots \\
1 & \omega_n^{-(n - 1)} & \omega_n^{-2(n - 1)} & \omega_n^{-3(n - 1)} & \dots & \omega_n^{-(n - 1)(n - 1)}
\end{bmatrix}
\begin{bmatrix}
g_0 \\
g_1 \\
\vdots \\
g_{n - 1}
\end{bmatrix}
\]
这也解释了 IDFT 的结论：将 DFT 中的 $\omega_n^1$ 换成 $\omega_n^{-1}$，做完以后除以 $n$。

\subsubsection{FFT 倍增实现}

\Highlight{位逆序置换}：

以 $n = 8$ 为例，模拟分治的过程：

\[
\begin{gathered}
\{ x_0, x_1, x_2, x_3, x_4, x_5, x_6, x_7 \} \\
\{ x_0, x_2, x_4, x_6 \}\{ x_1, x_3, x_5, x_7 \} \\
\{ x_0, x_4 \} \{ x_2, x_6 \} \{ x_1, x_5 \} \{ x_3, x_7 \} \\
\{ x_0 \} \{ x_4 \} \{ x_2 \} \{ x_6 \} \{ x_1 \} \{ x_5 \} \{ x_3 \} \{ x_7 \}
\end{gathered}
\]
从分治的角度，每次将 $x_i$ 按照奇偶性分类，偶则分到左边，奇则分到右边。直到只剩下一个项为止。

从倍增的角度，我们希望对于每个项 $x_i$，求出其经过不断分治以后，只剩下一个项时，在哪一个位置。然后根据变换后的位置，不断地倍增合并点值。

将 $x$ 用二进制数表示，变换后的位置即为\textbf{该二进制数的翻转}（例如 \texttt{011000} 变成 \texttt{000110}）。证明：每次根据奇偶性将 $x_i$ 进行分类，相当于是由最低位决定最高位。以此类推，次低位决定次高位，第三低位决定第三高位 ...

设 $\mathrm{rev}(x)$ 表示二进制数 $x$ 的翻转。首先有 $\mathrm{rev}(0) = 0$，考虑从小到大递推，有递推式

\[
\mathrm{rev}(x) = \left\lfloor \frac{\mathrm{rev}\left( \left\lfloor \frac{x}{2} \right\rfloor \right)}{2} \right\rfloor + (x \bmod 2) \times 2^{P - 1}
\]
\Highlight{蝶形运算}：已知 $F_L(\omega_{n / 2}^k), F_R(\omega_{n / 2}^k)$。每次更新，需要用 $F_L(\omega_{n / 2}^k) \pm \omega_n^kF_R(\omega_{n / 2}^k)$ 求出 $F(\omega_n^k), F(\omega_n^{k + n / 2})$。使用位逆序置换后

\begin{itemize}
\item 更新前：$F_L(\omega_{n / 2}^k)$ 储存在数组下标为 $k$ 的位置，$F_R(\omega_{n / 2}^k)$ 储存在数组下标为 $k + n / 2$ 的位置。
\item 更新后：$F(\omega_n^k)$ 储存在数组下标为 $k$ 的位置，$F(\omega_n^{k + n / 2})$ 储存在数组下标为 $k + n / 2$ 的位置。
\end{itemize}

因此可以直接在数组下标为 $k$ 和 $k + n / 2$ 的位置进行覆写，而不用开额外的数组保存值。

\subsubsection{FFT 多项式乘法}

对于两个多项式 $A, B$，现在要求出多项式 $C = A \times B$。

考虑对 $A, B$ 进行 DFT，经过 DFT 之后得到一系列关于 $A, B$ 的点值，将关于 $A, B$ 的点值对应位置相乘即可得到 $C$ 的点值，再对 $C$ 的点值进行 IDFT 即可得到 $C$ 的系数。

\begin{quoteblock}
核心思想：$A, B$ 的系数转 $A, B$ 的点值，$A, B$ 的点值点积转 $C$ 的点值，$C$ 的点值转 $C$ 的系数。
\end{quoteblock}

时间复杂度 $\mathcal{O}((n + m) \log (n + m))$。

\InputCode{cpp}{contents/0x30 数学/code/0x31 FFT.cpp}

\subsection{NTT}

回顾一些内容：

\Highlight{阶}：设整数 $m > 0$，满足 $a^x \equiv 1 \pmod m$ 的最小正整数 $x$，称作 $a$ 模 $m$ 的阶。记作 $\mathrm{ord}_m(a)$。

\Highlight{阶的存在条件}：\textbf{当且仅当 $\gcd(a, m) = 1$ 时}，$a$ 模 $m$ 的阶存在。

\Highlight{阶的简单性质}：

\begin{itemize}
\item 若 $a\perp m$，则 $a, a^2, \dots, a^{\mathrm{ord}_m(a)}$ 模 $m$ 两两不同。
\item 若 $a\perp m$，则 $\mathrm{ord}_m(a) \mid \varphi(m)$。
\item 若 $a \perp m$，则

\[
\mathrm{ord}_m(a^k) = \frac{\mathrm{ord}_m(a)}{\gcd(\mathrm{ord}_m(a), k)}
\]
\end{itemize}

\Highlight{原根}：设整数 $m > 0$，若 $\mathrm{ord}_m(g) = \varphi(m)$，则称 $g$ 为模 $m$ 的原根。

\Separator{}

在模意义下，我们希望找到一个东西来代替单位根。主要是要满足 FFT 过程中，用到的单位根满足的三个性质：

\begin{itemize}
\item $\omega_n$ 的阶为 $n$。
\item $\omega_{2n}^{2k} = \omega_n^k$。
\item $\omega_{2n}^k = -\omega_{2n}^{n + k}$。
\end{itemize}

对于一个质数 $m$，考虑 $m$ 的一个原根 $g$。此时 $g$ 的阶为 $m - 1$，显然原根 $g$ 不能直接代替 $n$ 阶单位根。

考虑 $g$ 的若干次幂。$g^k$ 的阶为 $\frac{m - 1}{\gcd(m - 1, k)}$，必为 $m - 1$ 的约数。\textbf{且当 $n \mid (m - 1)$ 时，可以构造 $g^{\frac{m - 1}{n}}$ 的阶恰好等于 $n$}。

尝试使用 $g^{\frac{m - 1}{n}}$ 代替 $n$ 阶单位根，性质 1, 2 的证明显然，性质 3 相当于是要证明 $g^{\frac{m - 1}{2}} \equiv -1 \pmod m$，可以使用二次剩余中的欧拉判别法来证明（原根必定不是二次剩余）。

\textbf{于是当 $n \mid (m - 1)$ 时，我们完全可以使用 $g^{\frac{m - 1}{n}}$ 来代替 $n$ 阶单位根 $\omega_n$。}

现在还要满足 $n \mid (m - 1)$ 才可以代替，对模数 $m$ 的取值有特殊的要求。注意到 FFT 过程中的 $n$ 都是 $2$ 的若干次幂，所以 $m - 1$ 只需要包含大量的因子 $2$ 即可。例如 $998244353 = 2^{23} \times 7 \times 17 + 1$，就是一个经典的 NTT 模数。

\Highlight{常用的 NTT 模数}：

\begin{itemize}
\item \texttt{998244353}：原根为 $3$。
\item \texttt{1004535809}：原根为 $3$。
\item \texttt{469762049}：原根为 $3$。
\item \texttt{167772161}：原根为 $3$。
\item \texttt{924844033}：原根为 $5$。
\item \texttt{4179340454199820289}（\textbf{$4 \times 10^{18}$ 级别，以时间换精度}）：原根为 $3$。
\end{itemize}

\InputCode{cpp}{contents/0x30 数学/code/0x31 NTT.cpp}

\subsection{MTT}

当模数不为常用的 NTT 模数时，我们就需要使用一些技巧使得多项式乘法可以继续进行。

首先不可以简单地使用 FFT，因为 double 的有效精度只有 15\textasciitilde{}16 位，而对于常见的数据范围 $n, m \leq 10^5$ 以及 $a_i, b_i \leq 10^9$，卷积结果的最大值接近 $10^5 \times 10^9 \times 10^9 = 10^{23}$，一定会掉精度。

\subsubsection{三模 NTT}

由于常见的数据范围下，卷积结果的最大值接近 $10^{23}$。于是我们考虑选择三个 $10^9$ 级别的常用 NTT 模数，\textbf{用三个模数分别做一次多项式乘法，将得到的三个结果使用 CRT 合并}。

常用的 NTT 模数选法如下，因为这三个 NTT 模数均有一个原根为 $3$。

\[
\begin{aligned}
& \mathrm{mod}1 = 998244353 \\
& \mathrm{mod}2 = 1004535809 \\
& \mathrm{mod}3 = 469762049
\end{aligned}
\]
CRT 过程中需要使用 \texttt{\_\_int128}。

\InputCode{cpp}{contents/0x30 数学/code/0x31 MTT（三模 NTT）.cpp}

\subsection{Trick：翻转卷积}

给出序列 $a_0, \dots, a_n$ 与序列 $b_0, \dots, b_m$，对于每一个 $k$ 需要求出 $c_k = \sum_{i} a_ib_{i - k}$，即下标差恒为 $k$ 的卷积式（相当于 $a_ib_j$ 贡献给了 $c_{i - j}$）。

将序列 $b$ 翻转得到 $b'$，其中 $b'_{m - j} = b_j$。\textbf{对序列 $a, b'$ 进行普通卷积得到序列 $c'$，此时 $a_ib_j$ 贡献给了 $c'_{m + i - j}$，于是有 $c_k = c'_{k + m}$}。

\subsection{Trick：积卷积}

给出序列 $a_1, \dots, a_{p - 1}$ 与序列 $b_1, \dots, b_{p - 1}$，对于每一个 $k \in [1, p)$ 需要求出 $c_k = \sum_{i\cdot j \bmod p = k} a_ib_j$。

当模数为较小的质数 $p$ 时，可以使用 “离散对数”，\textbf{将模 $p$ 意义下的乘法转换为模 $p - 1$ 意义下的加法}。

\subsection{Trick：多项式点乘向量}

给出多项式 $A(x) = \sum_{i = 0}^n a_ix^i$ 与向量 $\mathbf{B} = (b_0, \dots, b_n)$，现在我们要求多项式 $A$ 点乘向量 $\mathbf{B}$ 的值

\[
A \cdot \mathbf{B} = \sum_{i = 0}^n a_ib_i
\]
记翻转多项式 $B^R(x) = \sum_{i = 0}^n b_{n - i}x^i$，则点乘的值即为 $[x^n](A \times B^R)$，\textbf{即 $A \times B^R$ 的其中一项}。

当 $A$ 可以分解成两个多项式乘积 $A = A_1 \times A_2$ 时，若已知 $A_1 \times B^R$ 以及 $A_2$，可以 $\mathcal{O}(n)$ 求出 $[x^n](A \times B^R)$。

\begin{quoteblock}
e.g. 给出多项式序列 $p_1, \dots, p_n$，欲求 $\left(\prod_{i = l}^r p_i\right) \cdot \mathbf{B}$，即求 $\mathrm{pre}_r \times \mathrm{ipre}_{l - 1} \times B^R$。预处理出所有的 $\mathrm{pre}_i \times B^R$ 以及 $\mathrm{ipre}_i$ 即可。
\end{quoteblock}

\subsection{Model：分治卷积}

\Highlight{分治 NTT 模型 1}：给出 $n$ 个数 $a_1, \dots, a_n$，求出 $n$ 个单项式的积 $\prod_{i = 1}^n(a_ix + 1)$。

设 $f(l, r)$ 表示区间 $[l, r]$ 中所有单项式的积，\textbf{$f(l, r)$ 是一个次数为 $r - l + 1$ 的多项式（与区间长度有关）}。

取 $\mathrm{mid} = \lfloor \frac{l + r}{2} \rfloor$，则有

\[
f(l, r) = f(l, \mathrm{mid}) \times f(\mathrm{mid} + 1, r)
\]
这是一个分治的过程，不断分治直到 $l = r$ 即可。

时间复杂度 $T(n) = 2T\left( \frac{n}{2} \right) + \mathcal{O}(n \log n) = \mathcal{O}(n \log^2 n)$。

\begin{quoteblock}
By the way：求出 $n$ 个多项式的积，次数之和为 $m$。

可以每次取出两个次数最小的多项式相乘后放回，时间复杂度 $\mathcal{O}(m \log^2 m)$。其实使用分治求解的时间复杂度也是相同的。
\end{quoteblock}

\Highlight{分治 NTT 模型 2}：给出 $n$ 个数 $a_1, \dots, a_n$，求出 $n$ 个单项分式的和 $\sum_{i = 1}^n \frac{1}{1 - a_ix}$。

对于分式 $\frac{A_1}{A_2}, \frac{B_1}{B_2}$，若我们已知分子 $A_1, B_1$ 与分母 $A_2, B_2$，\textbf{由通分可得} $\frac{A_1}{A_2} + \frac{B_1}{B_2} = \frac{A_1B_2 + A_2B_1}{A_2B_2}$，由此可以得到新的分子 $A_1B_2 + A_2B_1$ 与分母 $A_2B_2$。

设 $A(l, r), B(l, r)$ 分别表示区间 $[l, r]$ 中所有单项分式的和，按照上述合并方法所得的分子与分母。$A(l, r), B(l, r)$ 的次数同样与区间长度有关。

取 $\mathrm{mid} = \lfloor \frac{l + r}{2} \rfloor$，则有

\[
\begin{gathered}
A(l, r) = A(l, \mathrm{mid}) \times B(\mathrm{mid} + 1, r) + A(\mathrm{mid} + 1, r) \times B(l, \mathrm{mid}) \\
B(l, r) = B(l, \mathrm{mid}) \times B(\mathrm{mid} + 1, r)
\end{gathered}
\]
通过分治求出 $A(1, n), B(1, n)$，最后求逆一次即可。

时间复杂度 $\mathcal{O}(n \log^2 n)$。

\Highlight{分治 NTT 模型 3}：给出 $n$ 个数 $a_1, \dots, a_n$，求出 $\sum_{i = 1}^n \prod_{j \neq i} (a_jx + 1)$。

设 $f(l, r) = \sum_{i = l}^r \prod_{j \in [l, r], j \neq i} (a_jx + 1)$ 与 $g(l, r) = \prod_{i = l}^r (a_ix + 1)$，$f(l, r), g(l, r)$ 的次数同样与区间长度有关。

取 $\mathrm{mid} = \lfloor \frac{l + r}{2} \rfloor$，则有

\[
\begin{gathered}
f(l, r) = f(l, \mathrm{mid}) \times g(\mathrm{mid} + 1, r) + f(\mathrm{mid} + 1, r) \times g(l, \mathrm{mid}) \\
g(l, r) = g(l, \mathrm{mid}) \times g(\mathrm{mid} + 1, r)
\end{gathered}
\]
时间复杂度 $\mathcal{O}(n \log^2 n)$。

\Highlight{分治 NTT 模型 4}：若递归到分治树的节点 $u$ 时，只需要保留多项式中的 $\mathcal{O}(s_u)$ 个有效项（其余的项可以特殊处理从而无需继续递归），并且节点 $u$ 上需要进行规模为 $\mathcal{O}(s_u)$ 的卷积。

记整棵分治树所有节点的 $s_u$ 之和等于 $m$，则时间复杂度 $\mathcal{O}(m \log m)$。

\Highlight{分治 NTT 模型 5（djq 分治）}：$Q$ 次询问组合数前缀和。每次询问给出 $n, m$，求 $\sum_{i = 0}^m \binom{n}{i} \bmod 998244353$。

注意到 $\binom{n}{m} = \frac{n \times (n - 1) \times \dots \times (n - m + 1)}{m!}$，可以将 $\binom{n}{m}$ 看成一个关于 $n$ 的 $m$ 次多项式。

记多项式 $F_i(x) = \frac{x - i + 1}{i}$（特别地 $F_0(x) = 1$），则单次查询的答案为

\[
\sum_{i = 0}^m\prod_{j = 0}^i F_j(n)
\]
记 $M = \max\{m\}$，对下标区间 $[0, M + 1]$ 建立线段树。对于一次询问 $n, m$，将其挂在叶子节点 $[m + 1, m + 1]$ 上。

对于线段树上表示区间 $[l, r]$ 的节点 $u$，预处理出：

\begin{itemize}
\item $\prod_{i = l}^r F_i(x)$。
\item $\sum_{i = l}^r \prod_{j = l}^i F_j(x)$。
\end{itemize}

接下来在线段树上递归，假设当前递归到了表示区间 $[l, r]$ 的节点 $u$，维护当前的 $\prod_{i = 0}^{l - 1} F_i(x)$ 与 $\sum_{i = 0}^{l - 1} \prod_{j = 0}^i F_j(x)$。在递归的过程中，两个多项式的次数太大，无法承受。

考虑运用多项式多点求值的经典套路：$f(x) \equiv f(n) \pmod{x - n}$。所以递归到节点 $u$ 的时候，将两个多项式均对 $\prod_{m + 1\in [l, r]} (x - n)$ 取模（也就是对区间内所有询问的 $x - n$ 之积取模，同样在线段树上预处理出，注意叶子可能有多个询问，乘积需要单独计算），此时两个多项式的次数均不超过区间内询问个数。

最后递归到叶子时，由于叶子可能有多个询问，对当前多项式进行多点求值即可。

时间复杂度 $\mathcal{O}(M \log^2 M + Q \log^2 Q)$，该做法的优点是时间复杂度与 $n$ 无关。

\begin{quoteblock}
还有一种基于莫队的做法。
\end{quoteblock}

\subsection{Model：半在线卷积}

\Highlight{半在线卷积}：给出 $g_1, \dots, g_n$，求出 $f_0, \dots, f_n$，其中初值 $f_0 = 1$，递推式 $f_i = \sum_{j = 1}^i f_{i - j}g_j$。

使用\textbf{强制中序遍历转移的 CDQ 分治}：设 $\mathrm{solve}(l, r)$ 表示求出区间 $[l, r]$ 中所有的 $f$ 值，取 $\mathrm{mid} = \lfloor \frac{l + r}{2} \rfloor$。

\begin{itemize}
\item 先递归 $\mathrm{solve}(l, \mathrm{mid})$，求出 $f_l, \dots, f_{\mathrm{mid}}$。
\item \textbf{计算 $f_{[l, \mathrm{mid}]}$ 与 $g_{[0, r - l]}$ 的卷积结果，求出左半区间对右半区间的贡献}。
\item 后递归 $\mathrm{solve}(\mathrm{mid} + 1, r)$，求出 $f_{\mathrm{mid} + 1}, \dots, f_r$。
\item 最后递归到 $l = r$ 时，请注意根据具体递推式调整 $f_i$（e.g. 递推式 $f_i = g_i - \sum_{j = 1}^{i - 1} \binom{i - 1}{j - 1} f_{i - j}g_j$）。
\end{itemize}

时间复杂度 $T(n) = 2T\left( \frac{n}{2} \right) + \mathcal{O}(n \log n) = \mathcal{O}(n \log^2 n)$。

\begin{quoteblock}
卡常技巧：计算 $f_{[l, \mathrm{mid}]}$ 与 $g_{[0, r - l]}$ 的卷积结果时，\textbf{只需将补充界设定成 $\geq \mathrm{len}$ 即可（而非 $\geq \frac{3 \cdot \mathrm{len}}{2}$）}。

实际上我们要求的是 $\frac{\mathrm{len}}{2}$ 次多项式与 $\mathrm{len}$ 次多项式相乘的 $\frac{\mathrm{len}}{2} \sim \mathrm{len}$ 次项。

由于 NTT 是循环卷积（$a_ib_j \to c_{(i + j) \bmod L}$，其中 $L$ 表示补充界），普通卷积的 $\mathrm{len} \sim \frac{3 \cdot \mathrm{len}}{2}$ 次项会污染掉循环卷积的 $0 \sim \frac{\mathrm{len}}{2}$ 次项，但我们只需要循环卷积的 $\frac{\mathrm{len}}{2} \sim \mathrm{len}$ 次项，所以将补充界调小不会对贡献产生影响。
\end{quoteblock}

\Highlight{全在线卷积}：求出 $f_0, \dots, f_n$，其中初值 $f_0 = 1$，递推式 $f_i = \sum_{j = 1}^i f_{i - j}g_j$，其中 $g_i$ 与 $f_0, \dots, f_i$ 强相关。

仍然使用强制中序遍历转移的 CDQ 分治。但 $g_{[0, r - l]}$ 我们不一定提前知道，进一步分析分治树的结构，\textbf{发现仅有 $l = 0$ 时不知道 $g_{[0, r - l]}$}。所以我们需要修改计算贡献的方式。

\begin{itemize}
\item 当 $l = 0$ 时，\textbf{计算 $f_{[0, \mathrm{mid}]} \times g_{[0, \mathrm{mid}]}$ 的贡献}（此时还有 $f_{[0, \mathrm{mid}]} \times g_{[\mathrm{mid} + 1, r]}$ 的贡献没被计算）。
\item 当 $l > 0$ 时，\textbf{计算 $f_{[l, \mathrm{mid}]} \times g_{[0, r - l]} + g_{[l, \mathrm{mid}]} \times f_{[0, r - l]}$ 的贡献}（后半部分即为 $l = 0$ 中漏掉的贡献）。
\end{itemize}

可以证明，点对的贡献不重不漏。

时间复杂度 $\mathcal{O}(n \log^2 n)$。

\subsection{Model：树上路径长度统计}

\Highlight{树上路径长度统计}：有一个包含 $n$ 个点的树，对于每个 $k$ ($1 \leq k < n$) 求出有多少条路径的长度为 $k$。

对树进行点分治，设当前的分治重心为 $c$，考虑统计经过当前分治重心的路径贡献。设删除分治重心 $c$ 之后，剩下的连通块内的点到 $c$ 的距离分布多项式分别为 $b_1(x), \dots, b_m(x)$，其中 $[x^k]b_i(x)$ 表示第 $i$ 个连通块有多少个距离 $c$ 为 $k$ 的点。

则分治重心 $c$ 的贡献多项式为

\[
\sum_{i = 1}^{m} b_i(x) \left( \sum_{j = 1}^{i - 1}b_j(x) + 1 \right)
\]
将多项式 $b_1(x), \dots, b_m(x)$ 按照\textbf{次数从小到大排序}后，维护一个前缀多项式之和，计算上式即可。

时间复杂度 $\mathcal{O}(n \log^2 n)$。

\begin{quoteblock}
还有一种基于长链剖分的做法，时间复杂度 $\mathcal{O}(n \log n)$，详见 \href{https://www.cnblogs.com/zkyJuruo/p/18809573}{zak 的博客}。
\end{quoteblock}

\subsection{多项式除法（暴力）}

\Highlight{多项式除法}：给出两个多项式 $A(x), B(x)$，你需要求出多项式 $C(x)$，满足 $B(x)C(x) = A(x)$。

设 $A(x) = \sum_{i = 0}^n a_ix^i, B(x) = \sum_{i = 0}^m b_ix^i, C(x) = \sum_{i = 0}^{n - m}c_ix^i$。直接根据定义 $\sum_{j = 0}^ic_jb_{i - j} = a_i$ 暴力展开。

先对 $b$ 的前导零进行处理，设 $b$ 的前导零有 $t$ 位，则需要使得 $a, b$ 左移 $t$ 位。

初值

\[
c_0 = \frac{a_0}{b_0}
\]
递推式

\[
c_i = \frac{1}{b_0}\left( a_i - \sum_{j = 0}^{i - 1}c_jb_{i - j} \right)
\]
但递推式里的 $j$ 只需要枚举到 $i - j \leq m$ 即 $j \geq i - m$ 就可以了。

时间复杂度 $\mathcal{O}(nm)$。实际上是 $(n - m)m$。

\subsection{多项式导数/积分}

\Highlight{多项式形式导数}：

\[
f'(x) = \left( \sum\limits_{i = 0}^na_ix^i \right)' = \sum\limits_{i = 1}^n ia_ix^{i - 1}
\]
\Highlight{多项式形式积分}：

\[
\int f(x) \mathrm{dx} = \int \left( \sum\limits_{i = 0}^n a_ix^i \right) \mathrm{dx} = \sum\limits_{i = 0}^n \frac{1}{i + 1}a_ix^{i + 1}
\]

\subsection{多项式牛顿迭代}

问题：已知函数 $g$，假设满足 $g(f(x)) = 0$ 的函数 $f(x) = h(x)$，我们需要求出 $h \bmod x^n$（在这里，我们将 $f(x)$ 看成函数 $g(f(x))$ 的自变量，则我们要求的就是 $g(f(x)) = 0$ 模 $x^n$ 意义下的零点）。

当 $n = 1$ 时，$[x^0]h(x)$ 需要单独求出。

当 $n > 1$ 时，\textbf{记 $h_0 = h \bmod x^{\left\lceil \frac{n}{2} \right\rceil}, h_1 = h \bmod x^n$，假设我们现在已经求出了 $h_0$，我们希望通过 $h_0$ 推出 $h_1$}。考虑将函数 $g(f(x))$ 在 $f(x) = h_0(x)$ 处泰勒展开。

\[
g(f(x)) = \sum\limits_{i = 0}^{\infty} \frac{g^{(i)}(h_0(x))}{i!}(f(x) - h_0(x))^i
\]
将 $h_1(x)$ 代入自变量 $f(x)$，观察到 $h_1(x) - h_0(x)$ 的前 $\left\lceil \frac{n}{2} \right\rceil$ 项均为 $0$（因为模 $x^{\left\lceil \frac{n}{2} \right\rceil}$ 意义下两者显然相等）。

也就是说 $h_1(x) - h_0(x)$ 的\textbf{最低次数}为 $\left\lceil \frac{n}{2} \right\rceil$，进一步有 $\left( h_1(x) - h_0(x) \right)^i$ 的最低次数为 $i\left\lceil \frac{n}{2} \right\rceil$。当 $i \geq 2$ 时，$\left(h_1(x) - h_0(x)\right)^i \bmod x^n$ 均为 $0$，\textbf{于是可以忽略 $i \geq 2$ 的项}。

进一步可以得到

\[
g(h_0(x)) + g'(h_0(x))\left( h_1(x) - h_0(x) \right) \equiv 0 \pmod{x^n}
\]
移项可得

\[
h_1(x) \equiv h_0(x) - \frac{g(h_0(x))}{g'(h_0(x))} \pmod{x^n}
\]
时间复杂度 $T(n) = T\left(\frac{n}{2}\right) + \mathcal{O}(n \log n) = \mathcal{O}(n \log n)$。

\subsection{多项式求逆}

\Highlight{多项式逆}：给出一个多项式 $A(x)$，你需要求出多项式 $B(x)$，满足 $A(x) \times B(x) = 1$。一般我们只需要求出 $B \bmod{x^m}$ 的结果。

\Highlight{多项式逆存在条件}：多项式 $A(x)$ 存在逆，当且仅当 $[x^0]A(x) \neq 0$。

（注意：\textbf{一个多项式的逆可能是有无穷项的}，例如 $\frac{1}{1 - x} = 1 + x + x^2 + \dots$）

\Highlight{少项式求逆（暴力）}：

设 $A(x) = \sum_{i = 0}^n a_ix^i, B(x) = \sum b_ix^i$。直接根据多项式逆的定义 $\sum_{j = 0}^{i} b_ja_{i - j} = [i = 0]$ 暴力展开。

初值

\[
b_0 = \frac{1}{a_0}
\]
递推式

\[
b_i = -\frac{1}{a_0} \sum\limits_{j = 0}^{i - 1}b_ja_{i - j}
\]
但递推式里的 $j$ 只需要枚举到 $i - j \leq n$ 即 $j \geq i - n$ 就可以了。

时间复杂度 $\mathcal{O}(nm)$。\textbf{当次数 $n$ 较小时有奇效}。

\begin{quoteblock}
这也验证了 “多项式逆存在条件” 的充分性。
\end{quoteblock}

\Highlight{多项式求逆（牛顿迭代）}：

构造函数 $g(f(x)) = \frac{1}{f(x)} - A(x)$，代入牛顿迭代法得

\[
\begin{gathered}
h_1(x) \equiv h_0(x) - \frac{\frac{1}{h_0(x)} - A(x)}{-\frac{1}{h_0^2(x)}} \pmod{x^m} \\
h_1(x) \equiv 2h_0(x) - A(x)h_0^2(x) \pmod{x^m}
\end{gathered}
\]

初始时，$A^{-1}(x)$ 的常数项即为 $A(x)$ 常数项的逆元。

时间复杂度 $\mathcal{O}(m \log m)$。

\subsection{多项式开方}

\Highlight{多项式开方}：给出一个多项式 $A(x)$，你需要求出多项式 $B(x)$，满足 $B^2(x) \equiv A(x)$。一般我们只需要求出 $B \bmod x^m$ 的结果。

\Highlight{多项式开方存在条件}：当 $[x^0]A(x) \neq 0$ 时，多项式 $A(x)$ 存在开方，当且仅当 $[x^0]A(x)$ 存在模意义下的二次剩余。

构造函数 $g(f(x)) = f^2(x) - A(x)$，代入牛顿迭代法得

\[
\begin{gathered}
h_1(x) \equiv h_0(x) - \frac{h_0^2(x) - A(x)}{2h_0(x)} \pmod{x^m} \\
h_1(x) \equiv \frac{h_0^2(x) + A(x)}{2h_0(x)} \pmod{x^m}
\end{gathered}
\]
初始时，$\sqrt{A(x)}$ 的常数项即为 $A(x)$ 常数项的二次剩余。当 $A(x)$ 常数项不等于 $1$ 时，还需要通过二次剩余求解常数项。

时间复杂度 $\mathcal{O}(m \log m)$。

\begin{quoteblock}
By the way：注意到 $\sqrt{A(x)} = e^{\frac{1}{2}\ln A(x)}$，可以直接套用多项式 ln 与多项式 exp 的模板。
\end{quoteblock}

\subsection{多项式带余除法}

\Highlight{多项式带余除法}：给出一个 $n$ 次多项式 $A(x)$ 和一个 $m$ 次多项式 $B(x)$，你需要求出多项式 $Q(x), R(x)$ 满足：

\begin{itemize}
\item $Q(x)$ 的次数为 $n - m$，$R(x)$ 的次数 $<m$。
\item $A(x) = Q(x) \times B(x) + R(x)$。
\end{itemize}

考虑构造变换 $f^R(x) = x^{\mathrm{deg} f} f\left(\frac{1}{x}\right)$，实际上即为翻转 $f$ 的系数。

将 $A(x) = Q(x) \times B(x) + R(x)$ 中的 $x$ 替换为 $\frac{1}{x}$，并且等式两边同乘以 $x^n$

\[
\begin{gathered}
x^nA\left(\frac{1}{x}\right) = x^{n - m}Q\left(\frac{1}{x}\right) \times x^mB\left(\frac{1}{x}\right) + x^{n - m + 1}x^{m - 1} R\left(\frac{1}{x}\right) \\
A^R(x) = Q^R(x) \times B^R(x) + x^{n - m + 1}x^{m - 1}R\left(\frac{1}{x}\right) \\
A^R(x) = Q^R(x) \times B^R(x) \pmod{x^{n - m + 1}}
\end{gathered}
\]
由于 $Q(x)$ 的次数小于 $n - m + 1$，故对 $x^{n - m + 1}$ 取模没有影响。

故在模 $x^{n - m + 1}$ 意义下求出 $\frac{A^R(x)}{B^R(x)}$ 再翻转一遍即可求出 $Q(x)$，再使用 $A(x) - Q(x) \times B(x)$ 即可求出 $R(x)$。

时间复杂度 $\mathcal{O}(n \log n)$。

\subsection{多项式 ln/exp}

\begin{quoteblock}
$e$ 在模 $998244353$ 下没有意义。同理，$\ln 2$ 在模 $998244353$ 下没有意义。

简单的解释：考察 $e = \sum_{i = 0}^\infty \frac{1}{i!}$，在 $i \geq 998244353$ 时，分式就已经失去了意义。即使有意义，无穷项求和在模 $998244353$ 下也不会收敛。
\end{quoteblock}

\Highlight{多项式 ln 与多项式 exp 的定义}：由麦克劳林级数定义

\[
\begin{aligned}
\ln f(x) & = \sum\limits_{i = 1}^{\infty} (-1)^{i + 1} \frac{(f(x) - 1)^i}{i} \\
& = - \sum\limits_{i = 1}^{\infty} \frac{(1 - f(x))^i}{i} \\
\exp f(x) & = \sum\limits_{i = 0}^{\infty} \frac{f^i(x)}{i!}
\end{aligned}
\]
\Highlight{多项式 ln 存在条件}：\textbf{多项式 $f(x)$ 存在 ln，当且仅当 $[x^0]f(x) = 1$，并且由此可得 $[x^0] \ln f(x) = 0$}。

\Highlight{多项式 ln}：给出一个多项式 $A(x)$，你需要求出多项式 $B(x) = \ln A(x)$。一般我们只需要求出 $B \bmod{x^m}$ 的结果。

将 $B(x) = \ln A(x)$ 等式两边对 $x$ 求导可得

\[
B'(x) = \frac{A'(x)}{A(x)}
\]
两边同时积分可得

\[
B(x) = \int \frac{A'(x)}{A(x)} \mathrm{dx}
\]
注意：虽然多项式先求导后积分会丢掉常数项的信息，但是容易验证当 $A(x)$ 的常数项为 $1$ 时，$B(x)$ 的常数项必定为 $0$。

时间复杂度 $\mathcal{O}(m \log m)$。

\Highlight{多项式 exp 存在条件}：\textbf{多项式 $f(x)$ 存在 exp，当且仅当 $[x^0]f(x) = 0$，并且可以推出 $[x^0]\exp f(x) = 1$}。

\Highlight{多项式 exp}：给出一个多项式 $A(x)$，你需要求出多项式 $B(x) = \exp A(x)$。一般我们只需要求出 $B \bmod{x^m}$ 的结果。

构造函数 $g(f(x)) = \ln f(x) - A(x)$，代入牛顿迭代得

\[
\begin{gathered}
h_1(x) \equiv h_0(x) - \frac{\ln h_0(x) - A(x)}{\frac{1}{h_0(x)}} \pmod{x^m} \\
h_1(x) \equiv h_0(x)\left( 1 - \ln h_0(x) + A(x) \right) \pmod{x^m}
\end{gathered}
\]
初始时，$\exp A(x)$ 的常数项必为 $1$。

时间复杂度 $\mathcal{O}(m \log m)$。

\subsection{多项式快速幂}

\Highlight{多项式快速幂}：给出一个多项式 $A(x)$，你需要求出 $A^k(x) \bmod{x^m}$。

\Highlight{多项式快速幂（倍增法）}：具体算法流程同快速幂。

时间复杂度 $\mathcal{O}(m \log m \log k)$，但可以处理的问题形式更加灵活。

\begin{quoteblock}
e.g. 多项式快速幂对另一多项式取模。

e.g. 多项式快速幂且保留次数在 $[-m, m]$ 的项。
\end{quoteblock}

\Highlight{多项式快速幂（先 ln 后 exp）}：注意到 $A^k(x) = e^{k\ln A(x)}$，考虑直接套用多项式 ln 与多项式 exp 的模板。

\textbf{但由于多项式 ln 的一些限制（需要保证 $[x^0]A(x) = 1$），首先要对 $A(x)$ 进行一些操作}。

找出一个最小的 $i$，满足 $[x^i]A(x) \neq 0$（相当于是找出 $A(x)$ 的第一个非零系数），记 $v = [x^i]A(x)$。则有

\[
A(x) = \left(\frac{A(x)}{x^iv}\right)x^iv
\]
\textbf{注意到 $\frac{A(x)}{x^iv}$ 满足常数项为 $1$ 的限制，可以进行多项式 ln}。

然后有

\[
A^k(x) = \left(\frac{A(x)}{x^iv}\right)^k x^{ik}v^k
\]
先使用多项式 ln 与多项式 exp 求出 $\left(\frac{A(x)}{x^iv}\right)^k \bmod{x^{m - ik}}$ 的值，最后再乘上 $x^{ik}v^k$ 即可。

（注意：这里的乘 $x^{ik}$ 相当于多项式系数向右移动 $ik$ 位。上面的除以 $x^i$ 同理，相当于多项式系数向左移动 $i$ 位，即舍弃掉前 $i$ 项系数）

时间复杂度 $\mathcal{O}(m \log m)$。

\begin{quoteblock}
设模数为质数 $P$。

\begin{itemize}
\item 当 $k \leq P$ 时，可以放心应用该做法。
\item 当 $k > P$ 时，\textbf{需要进一步讨论 $k$ 的取模相关问题}。
\end{itemize}

引理 1：设 $P$ 为质数。则对于任意整数 $a, b$

\[
(a + b)^P \equiv a^P + b^P \pmod P
\]

证明：因为 $(a + b)^P = \sum_{i = 0}^P \binom{P}{i}a^ib^{P - i}$，当 $0 < i < P$ 时，有 $\binom{P}{i} = \frac{P}{i}\binom{P - 1}{i - 1}$，由于 $P$ 是质数，所以该式子是 $P$ 的倍数。故上式得证。

引理 2：设 $f(x)$ 为一个 $n$ 次多项式，$P$ 为质数。则有

\[
f^P(x) \equiv f(x^P) \pmod P
\]

证明：不断地应用引理 1 即可。

由于一般情况下 $P > m$，所以 $f^P(x) \equiv f(x^P) \equiv ([x^0]f(x))^P \pmod{x^m}$。由于我们进行了一系列操作使得 $[x^0]f(x) = 1$，所以 $f^P(x) \equiv 1 \pmod{x^m}$。

\begin{itemize}
\item \textbf{所以多项式 ln 与多项式 exp 部分，$k$ 需要对 $P$ 取模}。
\item 对于 $x^{ik}$，当 $i = 0$ 时可以直接忽略，当 $i > 0$ 时 $k$ 显然不能超过 $P$，否则右移 $ik$ 位已经远超需要保留的范围。所以乘以 $x^{ik}$ 部分，当 $i > 0$ 时 $k$ 不能超过 $P$。
\item 对于 $v^k$，由于是数值计算，\textbf{所以乘以 $v^k$ 部分，$k$ 需要对 $\varphi(P) = P - 1$ 取模}。
\end{itemize}
\end{quoteblock}

\subsection{少项式快速幂}

\Highlight{少项式快速幂}：给出一个次数不高的多项式 $A(x)$，你需要求出 $A^k(x) \bmod x^m$。

设 $A(x) = \sum_{i = 0}^n a_ix^i, A^k(x) = \sum_{i = 0}^{m - 1}b_ix^i$。考虑对 $A^k(x)$ 求导

\[
\begin{gathered}
A^k(x)' = kA^{k - 1}(x)A'(x) \\
{\color{blue}A^k(x)'}A(x) = k{\color{blue}A^{k}(x)}A'(x)
\end{gathered}
\]
注意到 $A^k(x)$ 与 $A^k(x)'$ 有一定的递推关系。具体地，考虑等式左右两边的 $i$ 次项

\[
\sum\limits_{j = 0}^i (j + 1)b_{j + 1}a_{i - j} = k\cdot \sum_{j = 0}^i(i - j + 1)a_{i - j + 1}b_j
\]
等式左边只需枚举到 $j \geq i - n$，等式右边只需枚举到 $j \geq i - n + 1$。可以解得

\[
b_{i + 1} = \frac{\left( k\cdot \sum_{j = i - n + 1}^{i} (i - j + 1)a_{i - j + 1}b_j \right) - \left( \sum_{j = i - n}^{i - 1}(j + 1)b_{j + 1}a_{i - j} \right)}{(i + 1)a_0}
\]
初值 $b_0 = a_0^k$，上式可以在已知 $b_0, b_1, \dots, b_i$ 的情况下递推出 $b_{i + 1}$。

需要线性预处理出 $1, 2, \dots, m - 1$ 的逆元，递推式：$i^{-1} \equiv -\left\lfloor \frac{p}{i} \right\rfloor (p \bmod i)^{-1} \pmod{p}$。

时间复杂度 $\mathcal{O}(nm)$。\textbf{当次数 $n$ 较小时有奇效}。

\begin{quoteblock}
By the way：从这个递推式进行归纳可知，当 $a_0 = 1$ 时，$[x^n]A^k(x)$ 是一个关于 $k$ 的至多 $n$ 次多项式。
\end{quoteblock}

\subsection{*转置原理}

\sout{待填坑。}

\subsection{多项式多点求值}

\Highlight{多项式多点求值}：给出一个多项式 $f(x)$ 和 $n$ 个点 $x_1, x_2, \dots, x_n$，求 $f(x_1), f(x_2), \dots, f(x_n)$。

\Highlight{多项式多点求值（多项式取模）}：

考虑使用分治将问题规模减半。设 $\mathrm{mid} = \left\lfloor \frac{1 + n}{2} \right\rfloor$，将给定的点集分成两部分

\[
X_1 = \{ x_1,  \dots, x_{\mathrm{mid}} \}, \quad X_2 = \{ x_{\mathrm{mid} + 1}, \dots, x_{n} \}
\]
构造多项式

\[
g_1(x) = \prod\limits_{x_i \in X_1}(x - x_i), \quad g_2(x) = \prod\limits_{x_i \in X_2}(x - x_i)
\]
由多项式取模，记 $f(x) = Q_1(x) \times g_1(x) + R_1(x)$。\textbf{对于 $X_1$ 中的元素 $x_i$，由于 $g_1(x_i) = 0$，则有 $f(x_i) = R_1(x_i)$}。又由于 $R_1(x)$ 的次数小于区间长度的一半，于是问题规模成功减半。对于 $X_2$ 中的元素也是同理。

多项式 $g$ 可以通过分治 NTT 预处理。

时间复杂度 $T(n) = 2T(\frac{n}{2}) + \mathcal{O}(n \log n) = \mathcal{O}(n \log^2 n)$。

\begin{quoteblock}
另一种理解角度：$f(x) \equiv f(x_i) \pmod{x - x_i}$。证明只需将 $x^n = \left((x - x_i) + x_i\right)^n$ 二项式定理展开即可。
\end{quoteblock}

\InputCode{cpp}{contents/0x30 数学/code/0x31 多项式多点求值（多项式取模）.cpp}

\Highlight{多项式多点求值（转置原理）}：

\sout{待填坑。}

\subsection{多项式快速插值}

\Highlight{多项式快速插值}：给出 $n$ 个点 $(x_1, y_1), (x_2, y_2), \dots, (x_n, y_n)$，其中横坐标两两不同。快速地求一个至多为 $n - 1$ 次多项式，使得这 $n$ 个点都在 $f(x)$ 上。

考虑拉格朗日插值

\[
f(x) = \sum\limits_{i = 1}^{n}y_i \prod\limits_{j \neq i} \frac{x - x_j}{x_i - x_j}
\]
记 $g(x) = \prod_{j = 1}^n(x - x_j)$，则

\[
\prod_{j \neq i} (x_i - x_j) = \lim_{x \to x_i} \frac{\prod_{j = 1}^n (x - x_j)}{x - x_i} = g'(x_i)
\]
首先使用分治 NTT 求出 $g'(x)$，再使用多项式多点求值，求出所有的 $g'(x_i)$。

记 $v_i = \frac{y_i}{g'(x_i)}$，则原式化为

\[
f(x) = \sum\limits_{i = 1}^n v_i \prod\limits_{j \neq i} (x - x_j)
\]
仍然使用分治 NTT 求出上式。具体地，设 $\mathrm{mid} = \left\lfloor \frac{1 + n}{2} \right\rfloor$，设多项式

\[
\begin{gathered}
f_1(x) = \sum\limits_{i = 1}^\mathrm{mid} v_i \left( \prod\limits_{j \neq i, j \leq \mathrm{mid}} (x - x_j) \right), \quad g_1(x) = \prod\limits_{j \leq \mathrm{mid}} (x - x_j) \\
f_2(x) = \sum\limits_{i = \mathrm{mid} + 1}^n v_i \left( \prod\limits_{j \neq i, j > \mathrm{mid}} (x - x_j) \right), \quad g_2(x) = \prod\limits_{j > \mathrm{mid}} (x - x_j)
\end{gathered}
\]
显然有 $f(x) = f_1(x)g_2(x) + f_2(x)g_1(x)$。$f_1, f_2, g_1, g_2$ 可以进一步分治求解。

时间复杂度 $\mathcal{O}(n \log^2 n)$。

\subsection{多项式平移}

\Highlight{多项式平移}：给出一个多项式 $A(x)$ 以及一个常数 $c$，你需要求出多项式 $A(x + c)$。

\Highlight{多项式平移（二项式定理法）}：设 $A(x) = \sum_{i = 0}^n a_ix^i$，考虑二项式定理

\[
\begin{aligned}
A(x + c) & = \sum_{i = 0}^n a_i(x + c)^i \\
& = \sum_{i = 0}^n a_i \sum_{j = 0}^i \binom{i}{j}x^jc^{i - j} \\
& = \sum_{i = 0}^n a_i \sum_{j = 0}^i \frac{i!}{j!(i - j)!}x^jc^{i - j} \\
& = \sum_{i = 0}^n a_i \cdot i! \cdot \sum_{j = 0}^i \frac{x^j}{j!} \cdot \frac{c^{i - j}}{(i - j)!} \\
& = \sum_{j = 0}^n \frac{x^j}{j!} \cdot \left( \sum_{i = j}^n {\color{red} a_i \cdot i!} \cdot {\color{blue} \frac{c^{i - j}}{(i - j)!}} \right) \\
\end{aligned}
\]
注意到括号里的求和式，下标之差恒为 $j$，于是可以使用 “翻转卷积” 的技巧求出每个关于 $j$ 的求和式。

时间复杂度 $\mathcal{O}(n \log n)$。

\Highlight{连续点值平移}：给出 $n + 1$ 个连续点值 $f(0), f(1), \dots, f(n)$（其中 $f$ 是一个次数至多为 $n$ 的多项式）以及一个常数 $c$，你需要求出 $f(c), f(c + 1), \dots, f(c + n)$。

代入拉格朗日插值公式，得

\[
\begin{aligned}
f(x) & = \sum_{i = 0}^n y_i \prod_{j \neq i} \frac{x - j}{i - j} \\
& = \sum_{i = 0}^n y_i \cdot \frac{\prod_{v = x - n}^{x} v}{x - i} \cdot \frac{(-1)^{n - i}}{i!(n - i)!} \\
& = \left(\prod_{v = x - n}^{x} v\right) \sum_{i = 0}^n \frac{1}{x - i} \cdot \frac{y_i(-1)^{n - i}}{i!(n - i)!}
\end{aligned}
\]
于是

\[
f(c + x) = \left( \prod_{v = c + x - n}^{c + x} v \right)\sum_{i = 0}^n \frac{1}{c + x - i} \cdot \frac{y_i(-1)^{n - i}}{i!(n - i)!}
\]
注意到右边的求和式，是一个卷积的形式。\textbf{但是 $\frac{1}{c + x - i}$ 的分母在模意义下可能为 $0$}，此时 $c + x \equiv i \pmod{p}$，于是 $f(c + x) \equiv f(i) \equiv y_i \pmod{p}$，所以此时该点值不需要额外计算，直接取已有的值 $y_i$ 即可。

记：

\begin{itemize}
\item $a_i = \begin{cases} \frac{1}{c - n + i} & c - n + i \not\equiv 0 \pmod{p} \\ 0 & c - n + i \equiv 0 \pmod{p} \end{cases}$（其中 $i = 0 \sim 2n$，因为分母的取值范围为 $[c - n, c + n]$）。
\item $b_i = \frac{y_i(-1)^{n - i}}{i!(n - i)!}$。
\end{itemize}

则右边的求和式为 $\sum_{i = 0}^n a_{x - i + n}b_i$。将 $a, b$ 进行卷积得到 $\mathrm{res}$。此时，若 $(c + x) \bmod p \in [0, n]$，则

\[
f(c + x) = y_{(c + x) \bmod p}
\]
否则

\[
f(c + x) = \left( \prod_{v = c + x - n}^{c + x} v \right) \mathrm{res}_{x + n}
\]
时间复杂度 $\mathcal{O}(n \log n)$。

\begin{quoteblock}
给出 $n + 1$ 个等差点值 $f(b), f(b + k), \dots, f(b + nk)$（其中 $f$ 是一个次数至多为 $n$ 的多项式） 以及一个常数 $c$，你需要求出 $f(c + b), f(c + b + k), \dots, f(c + b + nk)$。

设 $g(x) = f(b + xk)$，可以看成：已知 $g(0), g(1), \dots, g(n)$，向右平移 $\frac{c}{k}$ 得到 $g(\frac{c}{k}), g(\frac{c}{k} + 1), \dots, g(\frac{c}{k} + n)$。
\end{quoteblock}

\subsection{多项式全家桶⭐}

% 待最终审稿：原 HTML <font size="6"> 开始
\textcolor{purple}{\textbf{详细的使用手册位于代码注释之中！}}
% 原 HTML </font> 结束

reference：\textcolor{black}{j}\textcolor{red}{iangly} 算法模板。

封装理念：

\begin{itemize}
\item \textbf{构建结构体 \texttt{poly} 继承自 \texttt{std::vector<int>}}，自由度高且易于访问。
\item 涉及 “当前多项式的变换或求值” 的功能均作为 \texttt{poly} 的内置函数；其余功能作为 \texttt{poly} 的外置函数。
\end{itemize}

\InputCode{cpp}{contents/0x30 数学/code/0x31 多项式全家桶.cpp}

\subsection{Bostan-Mori}

\subsubsection{分式远处系数}

用于求解分式 $\frac{P(x)}{Q(x)}$ 的第 $n$ 项 $[x^n]\frac{P(x)}{Q(x)}$。

考虑

\[
\frac{P(x)}{Q(x)} = \frac{P(x)Q(-x)}{Q(x)Q(-x)}
\]
记 $V(x) = Q(x)Q(-x)$，注意到 $V(x) = V(-x)$，则 $V(x)$ 仅有偶次项有值。于是我们有

\[
\frac{P(x)}{Q(x)} = \frac{C_{\mathrm{even}}(x^2)}{U(x^2)} + \frac{xC_{\mathrm{odd}}(x^2)}{U(x^2)}, \quad U(x^2) = Q(x)Q(-x)
\]
由于左边只影响偶次项，右边只影响奇次项（因为 $\frac{1}{U(x^2)} = \frac{1}{Q(x)Q(-x)}$ 也是偶函数，仅有偶次项有值），于是我们只需要根据 $n$ 的奇偶递归到其中一边即可。

当 $n$ 为偶数时

\[
[x^n]\frac{P(x)}{Q(x)} = [x^\frac{n}{2}]\frac{C_{\mathrm{even}}(x)}{U(x)}
\]
当 $n$ 为奇数时

\[
[x^n]\frac{P(x)}{Q(x)} = [x^{\frac{n - 1}{2}}]\frac{C_{\mathrm{odd}}(x)}{U(x)}
\]
递归到 $n = 0$ 时，有 $[x^0]\frac{f(x)}{g(x)} = \frac{f_0}{g_0}$。

时间复杂度 $\mathcal{O}(m \log m \log n)$，其中 $m$ 表示分子分母的长度。

\InputCode{cpp}{contents/0x30 数学/code/0x31 分式远处系数.cpp}

\subsubsection{分式远处区间系数}

\sout{待填坑。}

\subsection{常系数齐次线性递推}

\Highlight{常系数齐次线性递推}：已知初值 $a_0, \dots, a_{k - 1}$ 与系数 $c_1, \dots, c_k$（常数项 $c_0 = 0$），对于一个 $k$ 阶齐次线性递推数列 $a_i = \sum_{j = 1}^k c_ja_{i - j}$，你需要求出第 $n$ 项。

\subsubsection{Fiduccia 算法}

\Highlight{Fiduccia 算法}：构造特征多项式 $\Gamma(x) = x^k - \sum_{i = 1}^{k}c_ix^{k - i}$，记 $A(x) = \sum_{i = 0}^{k - 1} a_ix^i$，则有

\[
a_n = \left(x^n \bmod \Gamma(x) \right) \cdot A(x)
\]
（其中 $\cdot$ 代表点积）

于是应用多项式快速幂（倍增法）求出 $x^n \bmod \Gamma(x)$，最后求出点积即可。

时间复杂度 $\mathcal{O}(k \log k \log n)$，当项数 $n$ 与递推系数 $c$ 固定时，可以做到 $\mathcal{O}(k \log k \log n)$ 预处理，$\mathcal{O}(k)$ 查询。

\Highlight{Fiduccia 算法的简单解释}：我们想要知道 $a_n$ 具体可以被多少个 $a_0, \dots, a_{k - 1}$ 表示。相当于是初始有一个 $a_n$，每次选出最高项 $a_i$，根据公式 $a_i = \sum_{j = 1}^kc_ja_{i - j}$ 将 $a_i$ 消成更低的项。

从生成函数的角度考虑，相当于是将 $x^i$ 替换成 $\sum_{j = 1}^kc_jx^{i - j}$，等价于减去 $x^{i - k}\left( x^k - \sum_{j = 1}^kc_jx^{k - j} \right)$。\textbf{整个过程等价于 $x^n$ 对 $x^{k} - \sum_{j = 1}^k c_jx^{k - j}$ 取模}。

\InputCode{cpp}{contents/0x30 数学/code/0x31 常系数齐次线性递推（Fiduccia）.cpp}

\subsubsection{递推式转分式}

对于 $k$ 阶齐次线性递推数列 $\{a_i\}$，我们想要构造多项式 $P(x), Q(x)$ 使得 $[x^n]\frac{P(x)}{Q(x)} = a_n$。构造如下

\[
\begin{gathered}
Q(x) = 1 - \sum_{i = 1}^kc_ix^i \\
P(x) = \left(\left(\sum_{i = 0}^{k - 1} a_ix^i \right) \times Q(x)\right) \bmod{x^k}
\end{gathered}
\]
下面是检验：设 $\frac{P(x)}{Q(x)} = \sum_{i = 0}^{\infty} h_ix^i$。

\begin{itemize}
\item 当 $i = 0$ 时，有 $h_0 = \frac{p_0}{q_0} = a_0$。
\item 当 $1 \leq i < k$ 时，根据多项式 $P(x), Q(x)$ 的定义，显然有 $h_i = a_i$。
\item 当 $i \geq k$ 时，有 $h_i = -\frac{1}{q_0}\sum_{j = 1}^k q_jh_{i - j}$，化简得 $h_i = \sum_{j = 1}^k c_jh_{i - j}$，正是 $\{a_i\}$ 的递推式，归纳得 $h_i = a_i$。
\end{itemize}

\textbf{综上 $h_i = a_i$，原命题成立}。

于是构造多项式 $P(x), Q(x)$，套用 Bostan-Mori 求解 $[x^n]\frac{P(x)}{Q(x)}$ 即可。

时间复杂度 $\mathcal{O}(k \log k \log n)$，但单次询问常数相较于 Fiduccia 算法更小。

\InputCode{cpp}{contents/0x30 数学/code/0x31 常系数齐次线性递推（Bostan-Mori）.cpp}

\subsection{整式递推}

\sout{待填坑。}

\subsubsection{应用：快速阶乘算法}

\Highlight{快速阶乘算法}：求 $n! \bmod p$，其中 $p$ 是一个大质数。

（由于模数固定时，可以考虑分段打表。所以大概率不保证 $p$ 是一个 NTT 模数，此时需要使用 MTT）。

设 $v = \left\lfloor \sqrt{n} \right\rfloor$，设多项式 $g(x) = \prod_{i = 1}^v (x + i)$，则

\[
n! = \left( \prod_{i = 0}^{v - 1}g(iv) \right) \cdot \left( \prod_{i = v^2 + 1}^n i \right)
\]
首先后半部分的 $\prod_{i = v^2 + 1}^n i$ 可以 $\mathcal{O}(\sqrt{n})$ 求出。我们现在只关心前半部分。

设多项式 $g_d(x) = \prod_{i = 1}^d (x + i)$，则 $d + 1$ 个点值 $g_d(0), g_d(v), \dots, g_d(dv)$ 可以唯一确定 $d$ 次多项式 $g_d(x)$。

\textbf{所以对于当前的 $d$，我们都需要维护出这 $d + 1$ 个点值}。初始时 $d = 1$，只有点值 $g_1(0) = 1, g_1(v) = v + 1$。考虑两种操作：

操作 1：$d \gets 2d$。

\begin{itemize}
\item 补齐点值：补齐 $g_d((d + 1)v), \dots, g_d(2dv)$，可以通过 $g_d(0), g_d(v), \dots, g_d(dv)$ 向右平移 $dv$ 得到。\newline{}
\textcolor{blue}{（可以看作连续点值 $h(x) = g_d(xv)$ 向右平移 $d$）}
\item 修正点值：有公式 $g_{2d}(x) = g_d(x)g_d(x + d)$，将所有点值 $g_d(x)$ 向右平移 $d$ 即可得到 $g_d(x + d)$，然后相乘即可。\newline{}
\textcolor{blue}{（可以看作连续点值 $h(x) = g_d(xv)$ 向右平移 $\frac{d}{v}$）}
\end{itemize}

操作 2：$d \gets d + 1$。

\begin{itemize}
\item 修正点值：有公式 $g_{d + 1}(x) = g_d(x) \cdot (x + d + 1)$，直接相乘即可。
\item 补齐点值：补齐 $g_{d + 1}((d + 1)v)$，直接暴力 $\mathcal{O}(d)$ 求出即可。
\end{itemize}

于是可以从 $d = 1$ 通过倍增变换到 $d = v$。倍增过程的时间复杂度为 $T(n) = T(\frac{n}{2}) + \mathcal{O}(n \log n) = \mathcal{O}(n \log n)$。

而我们只需要 $\mathcal{O}(\sqrt{n})$ 个点值，故时间复杂度 $\mathcal{O}(\sqrt{n} \log n)$。

\InputCode{cpp}{contents/0x30 数学/code/0x31 快速阶乘算法.cpp}

\subsection{拉格朗日反演}

\sout{待填坑。}

